\newpage\handout
{Drill Problems 10}
{\textsc{Scholastic Aptitude Test (SAT)}}
{\href{https://creativecommons.org/licenses/by-nc-sa/4.0/}{CC BY-NC-SA 4.0 license}}
{Author: \BookAuthor}{Release: \generatedOn}
{SAT: Drill Problems (10)}
\input{drill00-disclaimer.tex}



\begin{enumerate}
  \item \textbf{Budget Interpretation} (10 points)\\
  Maximum Mail Order Company upgrades its computers each year. The function $P(t)=36 t+410$ models the company's budget for upgrades, in thousands of dollars, $t$ years after 2004. What is the best interpretation of 36 in this context?
  \begin{enumerate}[label=\Alph*.]
    \item The number of years until the budget is completely spent
    \item The amount budgeted, in thousands of dollars, in the year 2004
    \item The annual increase, in thousands of dollars, in the budget each year after 2004
    \item The number of years in which the budget will increase after 2004
  \end{enumerate}
  \begin{subanswer}
    % your answer here
  \end{subanswer}

  \item \textbf{Algebraic Expression Equivalence} (10 points)\\
  Which expression is equivalent to $\frac{4 x^2-3}{2 x+\sqrt{3}}$?
  \begin{enumerate}[label=\Alph*.]
    \item $2x-3$
    \item $2x-\sqrt{3}$
    \item $2x+\sqrt{3}$
    \item $2x+3$
  \end{enumerate}
  \begin{subanswer}
    % your answer here
  \end{subanswer}

  \item \textbf{$y$-intercept of Exponential Equation} (10 points)\\
  $$y-7=3^x-5$$\\
  What is the $y$-intercept for the given equation?
  \begin{subanswer}
    % your answer here
  \end{subanswer}

  \item \textbf{Linear Model for Calf Weight} (10 points)\\
  A linear model estimates that the mean weight of a calf in New Zealand, at birth, is 34 kilograms. At 180 days after birth, the mean weight is 272 kilograms. Which function best defines this model, where $f(x)$ is the mean weight of a calf, in kilograms, $x$ days after it is born, for $0 \leq x \leq 180$?
  \begin{enumerate}[label=\Alph*.]
    \item $f(x)=0.76 x+34$
    \item $f(x)=0.76 x+180$
    \item $f(x)=1.32 x+34$
    \item $f(x)=34 x+272$
  \end{enumerate}
  \begin{subanswer}
    % your answer here
  \end{subanswer}

  \newpage

  \item \textbf{Function Table Value} (10 points)\\
  \begin{table}[h!]
  \centering
  \renewcommand{\arraystretch}{1.3}
  \setlength{\tabcolsep}{16pt}
  \begin{tabular}{|c|c|}
  \hline $x$ & $f(x)$ \\
  \hline 1 & 3 \\
  \hline 3 & 6 \\
  \hline 5 & 9 \\
  \hline 7 & 12 \\
  \hline 9 & 15 \\
  \hline
  \end{tabular}
  \end{table}
  For the function $f$, the table shows some values of $x$ and their corresponding values of $f(x)$. If $f(1)=t$, what is the value of $f(3t)$?
  \begin{subanswer}
    % your answer here
  \end{subanswer}

  \item \textbf{Percent Increase from Chart} (10 points)\\
  \insertimage{0.40}{images/cus15}{reference attached}
  The chart shows the average cost of housing, in dollars per square foot, in Nebraska from 2015 to 2018. By approximately what percent did the average cost of housing increase over this period?
  \begin{enumerate}[label=\Alph*.]
    \item $20\%$
    \item $28\%$
    \item $39\%$
    \item $47\%$
  \end{enumerate}
  \begin{subanswer}
    % your answer here
  \end{subanswer}

  \newpage

  \item \textbf{Bacteria Growth Interpretation} (10 points)\\
  $$p(t)=2034(3)^{\frac{t}{14}}$$\\
  The given function shows the number of bacteria in a pond after $t$ days since the bacteria began to populate the pond. Which of the following is the best interpretation of $(3)^{\frac{t}{14}}$ in the equation?
  \begin{enumerate}[label=\Alph*.]
    \item The number of bacteria at the beginning of the study
    \item The number of bacteria triples every two weeks
    \item The number of bacteria increases by 3 every two weeks
    \item The number of bacteria in the pond after two weeks
  \end{enumerate}
  \begin{subanswer}
    % your answer here
  \end{subanswer}

  \item \textbf{Factoring Polynomial} (10 points)\\
  $$12 x^4-8 x^3-15 x^2$$\\
  The given expression is equal to $x^2(q x+5)(r x-3)$, where $q>0$ and $r>0$. What is the value of $q+r$?
  \begin{subanswer}
    % your answer here
  \end{subanswer}

  \item \textbf{Equation of Parallel Line} (10 points)\\
  Line $a$ is parallel to line $b$. Line $a$ goes through the points $(-4,0)$ and $(0,6)$. If line $b$ goes through the point $(-6,0)$, what is the equation for line $b$?
  \begin{enumerate}[label=\Alph*.]
    \item $y=\frac{2}{3} x-6$
    \item $y=\frac{2}{3} x+4$
    \item $y=\frac{3}{2} x+6$
    \item $y=\frac{3}{2} x+9$
  \end{enumerate}
  \begin{subanswer}
    % your answer here
  \end{subanswer}

  \item \textbf{Equilateral Triangle Exterior Angle} (10 points)\\
  The exterior angle of an equilateral triangle is equal to $(8x-20)^\circ$. What is the value of $x$?
  \begin{enumerate}[label=\Alph*.]
    \item 5
    \item 10
    \item 12.5
    \item 17.5
  \end{enumerate}
  \begin{subanswer}
    % your answer here
  \end{subanswer}

  \item \textbf{Quadratic Equation Solutions} (10 points)\\
  $$ (x-1)(x+2)=x(x-3)$$\\
  How many solutions does the given equation have?
  \begin{enumerate}[label=\Alph*.]
    \item 0
    \item 1
    \item 2
    \item 4
  \end{enumerate}
  \begin{subanswer}
    % your answer here
  \end{subanswer}

  \item \textbf{System of Inequalities Solution} (10 points)\\
  \[
    \begin{gathered}
    y>-2x-1 \\
    3y<x+9
    \end{gathered}
  \]
  Which of the following coordinates would be true for the given system of inequalities?
  \begin{enumerate}[label=\Alph*.]
    \item $(-3,3)$
    \item $(-2,1)$
    \item $(1,4)$
    \item $(3,1)$
  \end{enumerate}
  \begin{subanswer}
    % your answer here
  \end{subanswer}

  \item \textbf{Infinite Solutions in System of Equations} (10 points)\\
  \[
    \begin{aligned}
    9x-14y &= -3 \\
    ax-by &= 6
    \end{aligned}
  \]
  If the given system of equations has infinite solutions, what is the value of $b$?
  \begin{subanswer}
    % your answer here
  \end{subanswer}

  \newpage

  \item \textbf{Equivalent Forms of Quadratic Function} (10 points)\\
  $$f(x)=6x^2-45x+75$$\\
  Which of the following equivalent forms of the given function $f$ displays the $x$-intercepts as constants or coefficients?
  \begin{enumerate}[label=\Alph*.]
    \item $f(x)=(2x-5)(3x-15)$
    \item $f(x)=3x(2x-15)+75$
    \item $f(x)=6\left(x-\frac{15}{4}\right)^2+\frac{75}{8}$
    \item $f(x)=6(x-5)\left(x-\frac{5}{2}\right)$
  \end{enumerate}
  \begin{subanswer}
    % your answer here
  \end{subanswer}

  \item \textbf{System with One Real Solution} (10 points)\\
  \[
    \begin{gathered}
    y=4 \\
    y=x^2+12x+n
    \end{gathered}
  \]
  In the given system of equations, $n$ is a positive constant. For what value of $n$ does the system have exactly one real solution?
  \begin{subanswer}
    % your answer here
  \end{subanswer}

  \item \textbf{Minimum Value in System of Inequalities} (10 points)\\
  \[
    \begin{gathered}
    y \geq 2(x-3)^2+5 \\
    y \leq x+3
    \end{gathered}
  \]
  If $(a, b)$ is a solution to the given system of inequalities, what is the minimum value of $b$?
  \begin{enumerate}[label=\Alph*.]
    \item $-5$
    \item $-3$
    \item $3$
    \item $5$
  \end{enumerate}
  \begin{subanswer}
    % your answer here
  \end{subanswer}


  \newpage

  \item \textbf{Bacteria Doubling Problem} (10 points)\\
  Biology students grew a culture of bacteria at a constant temperature of $55^\circ\mathrm{C}$ for 3 days. The students began with 200 spores and the number of spores doubled every 12 hours. Of the equations given, which can be used to correctly calculate the number of spores in the bacteria culture at the end of 3 days?
  \begin{enumerate}[label=\Alph*.]
    \item $y=200(12)^3$
    \item $y=200(2)^3$
    \item $y=200(2)^6$
    \item $y=200(12)^2$
  \end{enumerate}
  \begin{subanswer}
    % your answer here
  \end{subanswer}

  \item \textbf{Comparing Means and Standard Deviations} (10 points)\\
  \insertimage{0.26}{images/cus16}{reference attached}
  Which of the statements best compares the means and standard deviations of Data Set $J$ and Data Set $K$?
  \begin{enumerate}[label=\Alph*.]
    \item The mean of Data Set $J$ is less than the mean of Data Set $K$, and the standard deviations of the two sets are equal.
    \item The mean of Data Set $J$ is greater than the mean of Data Set $K$, and the standard deviations of the two sets are equal.
    \item The means of the two sets are equal, and the standard deviation of Data Set $J$ is less than the standard deviation of Data Set $K$.
    \item The means of the two sets are equal, and the standard deviation of Data Set $J$ is greater than the standard deviation of Data Set $K$.
  \end{enumerate}
  \begin{subanswer}
    % your answer here
  \end{subanswer}

  \newpage

  \item \textbf{Circumference of a Circle} (10 points)\\
  What is the circumference of the circle with an equation of $x^2+6x+y^2-4y=51$?
  \begin{enumerate}[label=\Alph*.]
    \item $8\pi$
    \item $14.28\pi$
    \item $16\pi$
    \item $64\pi$
  \end{enumerate}
  \begin{subanswer}
    % your answer here
  \end{subanswer}

  \item \textbf{Trigonometric Value in Triangle} (10 points)\\
  \insertimage{0.40}{images/cus17}{reference attached}
  In the triangle shown, $\angle ABC=90^\circ$. Line segment $BD$ is shown with point $D$ lying on line $AC$. What is the value of $\sin \angle ABD-\cos \angle DBC$?
  \begin{enumerate}[label=\Alph*.]
    \item 0
    \item 30
    \item 45
    \item 90
  \end{enumerate}
  \begin{subanswer}
    % your answer here
  \end{subanswer}

  \item \textbf{No Solution for System of Equations} (10 points)\\
  \[
    \begin{gathered}
    y=(x-6)(x+6) \\
    y=k
    \end{gathered}
  \]
  In the system of equations shown, $k$ is an integer constant. What is the largest value of $k$ for which the system has no solution?
  \begin{subanswer}
    % your answer here
  \end{subanswer}

  \item \textbf{Exponential Function Shift} (10 points)\\
  The base of the exponential function $g(x)$ is $k$, where $k$ is a positive real number, and $k$ is equivalent to the value of the $y$-coordinate of the $y$-intercept of the graph of $g(x)$. Which of the following functions describes the exponential function $g(x-1)$?
  \begin{enumerate}[label=\Alph*.]
    \item $g(x-1)=k^x$
    \item $g(x-1)=k^{x+1}$
    \item $g(x-1)=k^x+k$
    \item $g(x-1)=k^{x-1}+k$
  \end{enumerate}
  \begin{subanswer}
    % your answer here
  \end{subanswer}

\end{enumerate}

